\documentclass[11pt,a4paper]{article}
\usepackage{brainteaser}

\begin{document}
\brainteasertitle{12}{Le pont et la lampe}

\begin{statementbox}
Quatre personnes doivent traverser un pont de nuit. Elles mettent
respectivement
\[
1,\quad2,\quad7,\quad10
\]
minutes pour traverser. Le pont supporte au maximum deux personnes; l'unique
lampe doit accompagner chaque traversée; et deux personnes avancent à la
vitesse de la plus lente.

Quel est le temps minimal permettant aux quatre personnes de traverser? Donner
une stratégie et prouver qu'elle est optimale.
\end{statementbox}

\begin{hintsbox}
\begin{enumerate}
  \item La stratégie où la personne de $1$ minute accompagne successivement
  les trois autres prend combien de temps?
  \item Comment faire traverser ensemble les personnes de $7$ et $10$ minutes?
  \item Qui doit rapporter la lampe avant et après leur traversée?
  \item Pour l'optimalité, séparer les cas où les deux plus lentes traversent
  ensemble ou séparément.
\end{enumerate}
\end{hintsbox}

\correction

Une stratégie en $17$ minutes est
\[
\begin{array}{rcl}
(1,2)&\longrightarrow&2,\\
1&\longleftarrow&1,\\
(7,10)&\longrightarrow&10,\\
2&\longleftarrow&2,\\
(1,2)&\longrightarrow&2.
\end{array}
\]
Le temps total vaut $2+1+10+2+2=17$.

Montrons qu'on ne peut pas faire mieux.

\subsection*{Les personnes de 7 et 10 minutes traversent séparément}
Les traversées vers l'arrivée coûtent au moins $10$, $7$ et $2$ minutes. Deux
retours sont nécessaires et coûtent au moins $1+1$. Le total est donc au moins
\[
10+7+2+1+1=21.
\]

\subsection*{Les personnes de 7 et 10 minutes traversent ensemble}
Leur traversée coûte $10$. Les deux passages nécessaires des personnes rapides
vers l'arrivée coûtent au moins $2+2$, et les deux retours nécessaires au
moins $1+2$. Le total est donc au moins
\[
10+4+3=17.
\]

\begin{resultbox}
La stratégie construite atteint la borne inférieure:
\[
\boxed{17\text{ minutes}.}
\]
\end{resultbox}

\end{document}
